\section{Introduction}

Neutrino-nucleus cross-sections have been the subject of intense study both experimentally and within the theory community in recent years due to their role in interpreting neutrino oscillation measurements and searches for other rare processes in neutrino scattering \cite{Benhar:2015wva}. While neutrino oscillation experiments primarily rely on measuring the rate of charged current (CC) interactions, it is also important that we build a solid understanding of inclusive and exclusive neutral current (NC) neutrino interactions.

NC neutrino interactions are of particular importance to $\nu_e$ and \nuebar{} measurements in the energy range of a few hundred MeV. This is especially true for detectors that cannot perfectly differentiate between photon- and electron-induced electromagnetic showers, and therefore where NC $\pi^0$ production followed by the subsequent decay $\pi^0 \rightarrow \gamma \gamma$ can be misidentified as $\nu_e$ or \nuebar{} CC scattering. Misidentification of photons as electrons complicates the interpretation of $\nu_e$ appearance measurements aiming to measure subtle signals. These include ~sterile neutrino oscillation searches with the upcoming Short Baseline Neutrino (SBN) experimental program \cite{SBNproposal} and CP violation measurements and mass hierarchy determination with the future Deep Underground Neutrino Experiment (DUNE) \cite{DUNE:2020ypp}. 

Furthermore, NC $\pi^0$ events can contribute as background to searches for rare neutrino scattering processes such as NC $\Delta$ resonance production followed by $\Delta$ radiative decay, or NC coherent single-photon production at energies below 1~GeV \cite{MicroBooNE:2021zai}. This is primarily a consequence of the limited geometric acceptance of some detectors, whereby one of the photons from a $\pi^0$ decay can escape the active volume of the detector. Depending on a detector's ability to resolve electromagnetic shower substructure, NC $\pi^0$ events can further contribute as background to searches for new physics beyond the Standard Model (BSM), such as $e^+e^-$ production predicted by a number of BSM models \cite{Bertuzzo:2018itn,Ballett:2018ynz,Abdullahi:2020nyr,Dutta:2020scq,Abdallah:2020vgg}.


Finally, NC measurements themselves can provide a unique channel for probing new physics. For example, searches for non-unitarity in the three-neutrino paradigm or searches for active to sterile neutrino oscillations are possible via NC rate disappearance measurements \cite{Cianci:2017okw,Furmanski:2020smg}. Such searches can provide complementary information to non-unitarity or light sterile neutrino oscillation parameters otherwise accessible only through CC measurements. 

Using a liquid argon time projection chamber (LArTPC) as its active detector, MicroBooNE \cite{MicroBooNE:2016pwy} shares the same technology and neutrino target nucleus as the upcoming SBN and future DUNE experiments. MicroBooNE's 85 metric ton active volume LArTPC is situated 468.5~m away from the proton beam target in the muon-neutrino-dominated Booster Neutrino Beam (BNB) at Fermilab~\cite{BNBflux} which is also used by SBN. The resulting neutrino beam has a mean energy $\langle E_\nu\rangle = 0.8$~GeV and is composed of $93.7\%~\nu_\mu$, $5.8\%~\bar{\nu}_\mu$, and $0.5\%~\nu_e/\bar{\nu}_e$. MicroBooNE's cross-section measurements on argon are therefore timely and directly relevant to these future (SBN and DUNE) programs.

We present the first measurement of neutrino-induced NC single-$\pi^0$ ($1\pi^0$) production on argon with a mean neutrino energy in the 1~GeV regime, which is also the highest-statistics measurement of this interaction channel on argon to date. This measurement is relevant to the physics programs of experiments that operate in the few-GeV regime (SBN~\cite{SBNproposal}, DUNE~\cite{DUNE:2020ypp}, NO$\nu$A~\cite{MINERvA:2016iqn,NOvA:2019bdw}, T2K~\cite{T2Kflux}, and Hyper-K~\cite{Hyper-KamiokandeProto-:2015xww}), especially those which share argon as a target material. Additionally, this measurement has been used to provide an indirect constraint to the rate of NC $1\pi^0$ backgrounds in MicroBooNE's recent search for a single-photon excess \cite{MicroBooNE:2021zai}. The only previous results for NC $1\pi^0$ scattering on argon are from the ArgoNeuT collaboration using the NuMI beam which has a much higher mean neutrino beam energy of 9.6 GeV for $\nu_\mu$ and of 3.6 GeV for $\overline{\nu}_\mu$  \cite{ArgoNeuT:2015ldo}. 

The interaction final states that are measured in this analysis are defined as
\begin{eqnarray}
\label{reaction}
\nu + \mathcal{A} \rightarrow \nu + \mathcal{A'} + \pi^0 + X,
\end{eqnarray}
where $\mathcal{A}$ represents the struck (argon) nucleus, $\mathcal{A'}$ represents the residual nucleus, and $X$ represents exactly one or zero protons plus any number of neutrons, but no other hadrons or leptons. The protons are identifiable in the MicroBooNE LArTPC by their distinct ionizing tracks while the $\pi^0$ is identifiable through the presence of two distinct electromagnetic showers, one for each photon from the $\pi^0\rightarrow\gamma\gamma$ decay, with kinematic properties such that they reconstruct to approximately the $\pi^0$ invariant mass.  

These one proton and zero proton samples are used first to perform a rate validation check and subsequently in three distinct cross-section measurements. By leveraging the capability of LArTPCs to detect and identify protons we perform the world's first exclusive NC$\pi^0$+0p and NC$\pi^0$+1p cross-section extractions and additionally measure the cross-section for NC$\pi^0$ interactions semi-inclusively using both the one proton and zero proton samples combined. Each of these cross-section extractions utilizes a distinct signal definition. The signal definitions for the two exclusive measurements place a threshold on true proton kinetic energy of greater than 50 MeV, while the semi-inclusive measurement allows for any number of protons. The signal definitions for all three measurements also require that there are no other hadrons or leptons in the final state (as noted above). MeV-scale photons, which may arise from nuclear de-excitation processes within the struck nucleus, are allowed in the final state. Finally, the signal definitions allow for interactions of all flavors of neutrinos that are present: $\nu_\mu$, $\bar{\nu}_\mu$, $\nu_e$, and $\bar{\nu}_e$.

These definitions are comparable to other historical NC $\pi^0$ measurements which typically require one and only one $\pi^0$ meson and little hadronic activity in the detector \cite{Barish:1974fe,Derrick:1980xw,Lee:1976wr,GargamelleNeutrinoPropane:1977hya,MiniBooNE:2008mmr,MiniBooNE2011,MiniBooNE:2009dxl,SciBooNE:2010NCpi0,K2K:2004qpv,T2K:2017epu}. This differs from the more inclusive approach of the ArgoNeuT experiment motivated both by its higher energy beam as well as the need to mitigate the low statistics of its data sample \cite{ArgoNeuT:2015ldo}. Making use of the MicroBooNE LArTPC's power in examining hadronic final state multiplicities and kinematic properties with high resolution, the flux-averaged cross-sections extracted in this analysis extend our understanding of this important interaction channel. The exclusive cross-sections reported provide new information useful for the tuning of NC $1\pi^0$ production and nuclear final state interactions in neutrino-argon scattering models, while the semi-inclusive cross-section enables a direct comparison to the MiniBooNE measurement of NC $\pi^0$ production.
